\documentclass[10pt,a4paper]{amsart}
\usepackage[margin=19mm]{geometry}
\usepackage{amsmath,amssymb,mathtools}
\usepackage[scr=rsfs,cal=euler]{mathalfa}
\usepackage{xfrac}
\usepackage[unicode,hidelinks,bookmarks=false]{hyperref}
\newcommand{\ie}{i.e.\ }
\newcommand{\cf}{cf.\ }
\newcommand{\resp}{respectively\ }
\newcommand{\Subsection}{\subsection}
\providecommand{\nobreakdash}{\nobreak\hbox{-}}
\setlength{\emergencystretch}{2em}
\multlinegap0pt
\usepackage{amssymb}
\usepackage{mathtools}
\usepackage{xcolor}
\newtheorem{lemma}{Lemma}
\newcommand{\AP}{\mathrm{AP}}
\newcommand{\F}{\mathbb F}
\title{Large Sets of Integers with No Harmonic Triples}
\author{Samuel Korsky}
\date{Source: arXiv:2607.05823v1 (CC BY 4.0)}
\begin{document}
\maketitle

\noindent\emph{Source and licence.} Large Sets of Integers with No Harmonic Triples. Version arXiv:2607.05823v1 (CC BY 4.0), distributed by arXiv under the Creative Commons Attribution 4.0 International licence, http://creativecommons.org/licenses/by/4.0/, as recorded in the arXiv licence metadata for this exact version and shown on the abstract page https://arxiv.org/abs/2607.05823v1. Full licence notice: \texttt{licenses/arxiv-2607.05823-CC-BY-4.0.txt}.\par\smallskip

\noindent\textit{Editorial scope.} Counted unit: the article's lemma lem:cardinality-core (source0.tex lines 170--177). The article's complete original proof of it is delivered with it (source0.tex lines 179--256, one \texttt{proof} environment of 78 lines). No mathematics is written, repaired or re-derived: the statement and the proof are the article's own lines, in the article's own notation, and no retained line is substituted or re-flowed.  Retained uncounted as the source-local context the proof needs: lines 65--67; lines 105--111; lines 127--141; lines 133--135.  Nothing was omitted from any retained span, so the package carries no omission record.  No retained line contains a citation, so the wrapper carries no bibliography and no bibliography entry.  No retained line contains a figure environment, an image-inclusion command or drawing code, so no image is delivered and the package declares no asset dependency.  Editorial formatting only, in addition: an AMS \texttt{amsart} preamble that restates the article's own declarations and macros used by the retained lines, namely the environments \texttt{lemma} on separate counters, together with the standard packages those lines need.  The retained environments print in this document as Lemma 1 in order of appearance, whereas the article prints the same items with the numbering of its own full document.  The complete native source of the article as distributed in the arXiv e-print for this exact version is bundled unchanged as \texttt{sources/204575/source0.tex}.
\begin{equation}\label{eq:harmonic-equation}
  \frac2a=\frac1b+\frac1c .
\end{equation}

\begin{lemma}[Progression-free subsets of prime fields]\label{lem:finite-field-input}
Let $q$ be prime and let $q\to\infty$.  There is a three-term-progression-free set $R\subseteq\F_q$ such that
\[
  \frac{|R|}{q}\ge \exp\!\left(-(C_{\AP}+o(1))\sqrt{\log q}\right),
  \qquad C_{\AP}=2\sqrt{\log(24/7)} .
\]
\end{lemma}

\begin{lemma}[Parametrization of odd harmonic triples]\label{lem:parametrization}
Let $b<a<c$ be odd positive integers satisfying
\begin{equation}\label{eq:param-equation}
  \frac2a=\frac1b+\frac1c .
\end{equation}
Then there exist positive integers $m,h,d$ with $1\le d<h$ such that
\begin{equation}\label{eq:parametrization}
  b=(h-d)hm,\qquad a=(h-d)(h+d)m,\qquad c=(h+d)hm .
\end{equation}
Conversely, every choice of positive integers $m,h,d$ with $1\le d<h$ gives a positive integer solution of \eqref{eq:param-equation} through \eqref{eq:parametrization}; it need not be an odd solution.  Moreover, if $q$ is an odd prime, $q\nmid abc$, and
\begin{equation}\label{eq:equal-congruence-param}
  a\equiv b\equiv c\pmod q,
\end{equation}
then $q\mid d$.
\end{lemma}

\begin{equation}
  b=(h-d)hm,\qquad a=(h-d)(h+d)m,\qquad c=(h+d)hm .
\end{equation}

\begin{lemma}[Cardinality construction]\label{lem:cardinality-core}
Let $Y$ be large.  There is a set $B\subseteq [Y]$ containing no distinct $a,b,c$ satisfying \eqref{eq:harmonic-equation} and such that
\begin{equation}\label{eq:core-cardinality}
  |B|
  \gg
  Y\exp\!\left(-(C_{\AP}+o(1))\sqrt{\log\log Y}\right).
\end{equation}
\end{lemma}

\begin{proof}
Let $K\ge 1$ be a sufficiently large absolute constant, to be fixed at the end of the proof.  By Bertrand's postulate, for all sufficiently large $Y$ there is a prime $q$ with
\begin{equation}\label{eq:q-choice}
  K\log Y\le q\le 2K\log Y .
\end{equation}
In particular, $q$ is odd and $\log q=\log\log Y+O(1)$.

Let $R\subseteq\F_q$ be supplied by Lemma~\ref{lem:finite-field-input}, and put
\[
  \delta=\frac{|R|}{q}.
\]
Thus
\begin{equation}\label{eq:delta-bound}
  \delta\ge \exp\!\left(-(C_{\AP}+o(1))\sqrt{\log\log Y}\right).
\end{equation}
For $\lambda\in\F_q^\times$ and $\mu\in\F_q$, write
\[
  R_{\lambda,\mu}=\lambda R+\mu .
\]
These affine images remain three-term-progression-free.

Let
\[
  V=\{n\le Y:n\text{ is odd and }q\nmid n\}.
\]
Then $|V|\gg Y$.  Choose $\lambda\in\F_q^\times$ and $\mu\in\F_q$ uniformly at random, and define
\[
  B_0=B_0(\lambda,\mu)=\{n\in V:n^{-1}\bmod q\in R_{\lambda,\mu}\}.
\]
For each fixed nonzero residue $r\in\F_q$, the probability that $r\in R_{\lambda,\mu}$ is exactly $\delta$.  Hence
\begin{equation}\label{eq:initial-expected-size}
  \mathbb E\!\left[|B_0|\right]=\delta |V|\gg \delta Y .
\end{equation}

We next delete the largest element of every surviving harmonic triple. For a fixed affine image $R_{\lambda,\mu}$, delete from $B_0$ every element $c$ that is the largest member of some triple $b<a<c$ in $B_0$ satisfying \eqref{eq:harmonic-equation}.  Call the remaining set $B=B(\lambda,\mu)$.  This set is harmonic-triple-free, since the largest element of any surviving harmonic triple would have been deleted.

It remains to bound the expected number of deletions.  Suppose $b<a<c$ is a harmonic triple in $B_0$.  Reducing \eqref{eq:harmonic-equation} modulo $q$ gives
\[
  2a^{-1}\equiv b^{-1}+c^{-1}\pmod q.
\]
The three residues $b^{-1},a^{-1},c^{-1}$ lie in the progression-free set $R_{\lambda,\mu}$, so they must all be equal.  Therefore
\begin{equation}\label{eq:abc-congruent}
  a\equiv b\equiv c\pmod q .
\end{equation}
Since $b,a,c\in V$, they are odd and not divisible by $q$.  By Lemma~\ref{lem:parametrization}, every such triple has the form \eqref{eq:parametrization} with $q\mid d$.

Let
\[
  U_q(Y)=
  \#\{(h,d,m):1\le d<h,\ q\mid d,\ h(h+d)m\le Y\}.
\]
This is an overcount for the possible triples whose largest element can be deleted.  If a single element $c$ is deleted because of several triples, then counting all of those triples only increases the upper bound.  We have
\begin{align}
  U_q(Y)
  &\le
  \sum_{h\le Y^{1/2}}
  \sum_{\substack{1\le d<h\\ q\mid d}}
  \left\lfloor \frac{Y}{h(h+d)}\right\rfloor \notag\\
  &\le
  Y\sum_{h\le Y^{1/2}}
  \sum_{1\le j<h/q}
  \frac1{h(h+qj)} \notag\\
  &\ll
  Y\sum_{h\le Y^{1/2}}\frac1{qh}
  \ll \frac{Y\log Y}{q}.\label{eq:Uq-bound}
\end{align}
For each parametrized triple counted by $U_q(Y)$, the probability that it contributes to a deletion is at most $\delta$.  Indeed, if the displayed integers are not all odd, or if one of them is divisible by $q$, then this probability is zero.  Otherwise $q\mid d$ implies that the three entries are congruent to the same nonzero residue modulo $q$, so their inverse residues are also equal; the probability that this common inverse residue lies in $R_{\lambda,\mu}$ is exactly $\delta$.  By the union bound, the expected number of deleted elements is at most $\delta U_q(Y)$.

Combining \eqref{eq:initial-expected-size}, \eqref{eq:Uq-bound}, and \eqref{eq:q-choice}, we get
\[
  \mathbb E\!\left[|B|\right]
  \ge
  \delta\left(c_1Y-c_2\cdot\frac{Y\log Y}{q}\right)
  \ge
  \delta\left(c_1Y-\frac{c_2}{K}\cdot Y\right)
\]
for absolute constants $c_1,c_2>0$.  Taking $K>2c_2/c_1$, the expectation is $\gg \delta Y$.  Therefore some affine image yields a set $B$ satisfying \eqref{eq:core-cardinality}, by \eqref{eq:delta-bound}.
\end{proof}

\end{document}
